% Inbuilt themes in beamer

\documentclass[aspectratio=169]{beamer}
\usepackage[export]{adjustbox}
\usepackage{xcolor}

\definecolor{darkred}{RGB}{176, 31, 36}
\definecolor{darkgreen}{RGB}{44, 95, 47}

\usepackage{listings}
\usepackage{booktabs}
\usepackage{graphicx}

% Title page details: 
\title{Tema 3.1: Redes recurrentes}
\subtitle{Fundamentos y arquitecturas}

% Theme choice:
\usetheme{Madrid}
\usecolortheme{seagull}
\useinnertheme{circles}

\author{Pablo Pérez Núñez}
 
\setbeamertemplate{navigation symbols}{}
\setbeamertemplate{author}{}

% Ajusta las dimensiones de la diapositiva

\setbeamertemplate{blocks}[rounded][shadow=false]
% \setbeamertemplate{blocks}[default]
\definecolor{myblue}{RGB}{27, 122, 194}
\definecolor{lightblue}{RGB}{55, 150, 225}
\definecolor{lightbluetwo}{RGB}{164, 209, 242}

\setbeamercolor{block title}{bg=lightblue,fg=white}
\setbeamercolor{block body}{bg=lightbluetwo,fg=black}

\setbeamercolor{title}{fg=white, bg =lightblue}
\setbeamercolor{frametitle}{fg=white, bg = lightblue}

\setbeamercolor*{palette secondary}{fg=white,bg=lightblue}

\defbeamertemplate*{footline}{infolines}
{
  \leavevmode%
  \hbox{%
  \begin{beamercolorbox}[wd=\paperwidth,ht=2.25ex,dp=1ex,left]{title in head/foot}%
    \hspace*{2ex} \usebeamerfont{title in head/foot}\insertshorttitle \ - \insertsubtitle
  \end{beamercolorbox}
  % \begin{beamercolorbox}[wd=.25\paperwidth,ht=2.25ex,dp=1ex,right]{date in head/foot}%
  %   \usebeamerfont{date in head/foot}\hspace*{2em}
  %   \insertframenumber{} / \inserttotalframenumber\hspace*{2ex} 
  %\end{beamercolorbox}
  }%
  \vskip0pt%
}

\setbeamertemplate{title page}{
  \centering
  % \usefonttheme{serif}
  \begin{minipage}{0.95\textwidth} % Ajustar el ancho según sea necesario

    \vspace{1em}
    {\color{myblue}\rule{\linewidth}{2pt}\par}
    \vspace{.5cm}
    {\bfseries \huge \textcolor{myblue}{\inserttitle}\par }
    \vspace{1em}
    {\bfseries \LARGE \textcolor{myblue}{\insertsubtitle}\par }
      
    \vspace{2.5em}

    \begin{columns}
      \column{0.6\textwidth}
      \vspace{6em}\\
      {\bfseries \large \textcolor{myblue}{Aprendizaje profundo / Deep Learning}} \\
      \vspace{1em}
      {\textcolor{myblue}{\insertauthor}}
      \column{0.4\textwidth}
      \vfill
      \centering
      \includegraphics[height=4.5em]{imgs/LogoAEPIA.pdf} \\
      \vspace{1em}
      \includegraphics[height=2.5em]{imgs/UIMP.png}
    \end{columns}

    \vspace{1em}

    {\color{myblue}\rule{\linewidth}{.5pt}\par}

  \end{minipage}
    

  %  \usefonttheme{default}
}

% Define the template for section slides
\AtBeginSection[]{
  \begin{frame}
    \vfill
    \LARGE\insertsectionhead\par
    {\color{myblue}\rule{\linewidth}{1pt}\par}
    \vfill
  \end{frame}
}
 

% MIS COSAS --------------------------------
% Bloque con margen inferior
\newenvironment{blockm}[1]{%
  \begin{block}{\textbf{#1}}%
  }{%
  \end{block}%
  \vspace{1em}%
}
% Dos imágenes
\newcommand{\twoimages}[3]{%
  \begin{columns}
    \begin{column}{0.5\textwidth}
      \centering
      \includegraphics[width=#3\textwidth]{#1}
    \end{column}
    \begin{column}{0.5\textwidth}
      \centering
      \includegraphics[width=#3\textwidth]{#2}
    \end{column}
  \end{columns}
}

% Comando para mostrar el título de la subsección al inicio de cada nueva subsección
\AtBeginSubsection[]
{
  \begin{frame}
    \vfill
    \small\insertsectionhead\par
    {\color{myblue}\rule{\linewidth}{1pt}\par}
	\LARGE\insertsubsectionhead\par
    \vfill
  \end{frame}
}

% ------------------------------------------

\begin{document}

% Title page frame
\begin{frame}[plain]
	\titlepage 
\end{frame}

\logo{}

\begin{frame}{Introducción}
  
  \begin{blockm}{}
    En la actualidad existen numerosas arquitecturas DNN creadas en función del problema que pretenden resolver.\\
  \end{blockm}
  Veremos en detalle las siguientes:
  \begin{itemize}
    \item Redes recurrentes
    \item Transformers
    \item Autoencoders
	\item Redes generativas adversarias
  \end{itemize}
\end{frame}

% ---------------------------------------------
\section{Redes neuronales recurrentes (RNN)}

\begin{frame}{Introducción}
  \begin{blockm}{¿Qué son?}
    Redes neuronales profundas que incorporan capas con conexiones recurrentes.
  \end{blockm}
  Las conexiones recurrentes sirven para sacar partido de lo aprendido en instantes anteriores.
  \vspace{1em}
  \twoimages{imgs/tema4/rnn/DNN.pdf}{imgs/tema4/rnn/RNN.pdf}{0.55}
\end{frame}

\begin{frame}{Introducción}
  \begin{blockm}{¿Para qué sirven?}
    Son útiles en el tratamiento de información secuencial (texto, audio, video,...).
  \end{blockm}
  Ejemplos:
  \begin{itemize}
    \item Reconocimiento del habla
    \item Análisis de sentimientos a partir de textos
    \item Traducción automática
    \item Reconocimiento de entidades
    \item Generación de música
    \item Análisis de secuencias de ADN
    \item etc.
  \end{itemize}

\end{frame}

\begin{frame}{RNN vs Feed Forward}
  \begin{blockm}{\textbf{Ejemplo:} Predecir la palabra siguiente a partir del comienzo de una frase}
    \textit{Entrada:} Secuencia de palabras de longitud variable.\\
    \textit{Salida:} Palabra.
  \end{blockm}
  Otros datos:
  \begin{itemize}
    \item Cada palabra se codifica utilizando one-hot.
    \item Una frase es una secuencia con número variable de vectores one-hot.
    \item Supongamos:
    \begin{itemize}
      \item Vocabulario de 10000 palabras.
      \item Longitud máxima de frase de 50 palabras.
    \end{itemize}
  \end{itemize}
\end{frame}

\begin{frame}{Feed Forward vs RNN}
  Mismo problema, arquitecturas diferentes:
  \vspace{1em}
  \twoimages{imgs/tema4/rnn/EX1.pdf}{imgs/tema4/rnn/EX2.pdf}{0.75}
\end{frame}


\begin{frame}{Notación}
	Matrices de pesos:
	\begin{itemize}
	  \item $\mathbf{U}$: Conexiones entre capa de entrada y recurrente.
	  \item $\mathbf{V}$: Conexiones de la capa recurrente con la siguiente capa.
	  \item $\mathbf{W}$: Conexiones de la capa recurrente con ella misma.
	  \item $\mathbf{b}_{h}, \mathbf{b}_{y}$: Bias de la capa recurrente y de la de salida.
	\end{itemize}
	\vspace{1em}
	Funciones de activación:
	\begin{itemize}
	  \item $g_{1}$: Suele ser Tangente Hiperbólica o ReLU.
	  \item $g_{2}$: Sigmoide (generalmente).
	\end{itemize}
  \end{frame}

\begin{frame}{Notación}

	\begin{columns}
		\begin{column}{0.4\textwidth}
			Dimensiones:
			\begin{itemize}
			  \item $\mathbf{U}$: $3x5$
			  \item $\mathbf{V}$: $2x3$
			  \item $\mathbf{W}$: $3x3$
			  \item $\mathbf{b}_{h}$: $3x1$
			  \item $\mathbf{b}_{y}$: $2x1$
			\end{itemize}
			\vspace{1em}
			Ecuaciones:
			\begin{itemize}
			  \item $\mathbf{h}^{(t)} = g_{1}(\mathbf{Wh}^{(t-1)}+\mathbf{Ux}^{(t)}+\mathbf{b}_{h})$
			  \item $\mathbf{\hat{y}}^{(t)} = g_{2}(\mathbf{Vh}^{(t)}+\mathbf{b}_{y})$
			\end{itemize}
		\end{column}

		\begin{column}{0.6\textwidth}
			%Detalle de los elementos de una RNN:\\
			\begin{center}
				\includegraphics[width=\textwidth]{imgs/tema4/rnn/ECU.pdf}
			\end{center}
		\end{column}
	\end{columns}

\end{frame}

\begin{frame}{Memoria en las RNN}

	\begin{columns}
		\begin{column}{0.5\textwidth}
			Partiendo de:
			\begin{itemize}
				\item $\mathbf{h}^{(t)} = g_{1}(\mathbf{Wh}^{(t-1)}+\mathbf{Ux}^{(t)}+\mathbf{b}_{h})$
			\end{itemize}			
			Para simplificar, eliminamos activación y bias:
			\begin{itemize}
				\item $\mathbf{h}^{(t)} = \mathbf{Wh}^{(t-1)}+\mathbf{Ux}^{(t)}$
			\end{itemize}
			\vspace{1em}
			Notación:
			\begin{itemize}
			  \item $\mathbf{x}^{(i)}$: Vector del i-ésimo elemento de la secuencia de entrada.
			  \item $\mathbf{h}^{(i)}$: Salida de la capa recurrente debida a la entrada $\mathbf{x}^{(i)}$, es decir, en la etapa i-ésima.
			\end{itemize}
		\end{column}

		\begin{column}{0.5\textwidth}
			%Detalle de los elementos de una RNN:\\
			\begin{center}
				\includegraphics[width=.8\textwidth]{imgs/tema4/rnn/Memoria1.pdf}
			\end{center}
		\end{column}
	\end{columns}

\end{frame}

\begin{frame}{Memoria en las RNN}

	\begin{columns}
		\begin{column}{0.5\textwidth}
			Partiendo de:
			\begin{itemize}
				\item $\mathbf{h}^{(t)} = g_{1}(\mathbf{Wh}^{(t-1)}+\mathbf{Ux}^{(t)}+\mathbf{b}_{h})$
			\end{itemize}			
			Para simplificar, eliminamos activación y bias:
			\begin{itemize}
				\item $\mathbf{h}^{(t)} = \mathbf{Wh}^{(t-1)}+\mathbf{Ux}^{(t)}$
			\end{itemize}
			\vspace{1em}
			Notación:
			\begin{itemize}
			  \item $\mathbf{x}^{(i)}$: Vector del i-ésimo elemento de la secuencia de entrada.
			  \item $\mathbf{h}^{(i)}$: Salida de la capa recurrente debida a la entrada $\mathbf{x}^{(i)}$, es decir, en la etapa i-ésima.
			\end{itemize}
		\end{column}

		\begin{column}{0.5\textwidth}
			%Detalle de los elementos de una RNN:\\
			\begin{center}
				\includegraphics[width=.8\textwidth]{imgs/tema4/rnn/Memoria2.pdf}
			\end{center}
		\end{column}
	\end{columns}

\end{frame}

\begin{frame}{Memoria en las RNN}

	\begin{columns}
		\begin{column}{0.5\textwidth}
			Partiendo de:
			\begin{itemize}
				\item $\mathbf{h}^{(t)} = g_{1}(\mathbf{Wh}^{(t-1)}+\mathbf{Ux}^{(t)}+\mathbf{b}_{h})$
			\end{itemize}			
			Para simplificar, eliminamos activación y bias:
			\begin{itemize}
				\item $\mathbf{h}^{(t)} = \mathbf{Wh}^{(t-1)}+\mathbf{Ux}^{(t)}$
			\end{itemize}
			\vspace{1em}
			Notación:
			\begin{itemize}
			  \item $\mathbf{x}^{(i)}$: Vector del i-ésimo elemento de la secuencia de entrada.
			  \item $\mathbf{h}^{(i)}$: Salida de la capa recurrente debida a la entrada $\mathbf{x}^{(i)}$, es decir, en la etapa i-ésima.
			\end{itemize}
		\end{column}

		\begin{column}{0.5\textwidth}
			%Detalle de los elementos de una RNN:\\
			\begin{center}
				\includegraphics[width=.8\textwidth]{imgs/tema4/rnn/Memoria3.pdf}
			\end{center}
		\end{column}
	\end{columns}

\end{frame}

\begin{frame}{Memoria en las RNN}

	\begin{columns}
		\begin{column}{0.5\textwidth}
			Partiendo de:
			\begin{itemize}
				\item $\mathbf{h}^{(t)} = g_{1}(\mathbf{Wh}^{(t-1)}+\mathbf{Ux}^{(t)}+\mathbf{b}_{h})$
			\end{itemize}			
			Para simplificar, eliminamos activación y bias:
			\begin{itemize}
				\item $\mathbf{h}^{(t)} = \mathbf{Wh}^{(t-1)}+\mathbf{Ux}^{(t)}$
			\end{itemize}
			\vspace{1em}
			Notación:
			\begin{itemize}
			  \item $\mathbf{x}^{(i)}$: Vector del i-ésimo elemento de la secuencia de entrada.
			  \item $\mathbf{h}^{(i)}$: Salida de la capa recurrente debida a la entrada $\mathbf{x}^{(i)}$, es decir, en la etapa i-ésima.
			\end{itemize}
		\end{column}

		\begin{column}{0.5\textwidth}
			%Detalle de los elementos de una RNN:\\
			\begin{center}
				\includegraphics[width=.8\textwidth]{imgs/tema4/rnn/Memoria4.pdf}
			\end{center}
		\end{column}
	\end{columns}

\end{frame}

\begin{frame}[shrink=25]{Memoria en las RNN}
	\begin{columns}
		\begin{column}{0.6\textwidth}
			{\color{darkred}
			\begin{equation*}
				\begin{aligned}
					& \mathbf{h}^{(1)}=g_{1}\left(\mathbf{W h}^{(0)}+\mathbf{U} \mathbf{x}^{(1)}+\mathbf{b}_{h}\right) \\
					& \mathbf{h}^{(2)}=g_{1}\left(\mathbf{W}\left(g_{1}\left(\mathbf{W} \mathbf{h}^{(0)}+\mathbf{U} \mathbf{x}^{(1)}+\mathbf{b}_{h}\right)\right)+\mathbf{U} \mathbf{x}^{(2)}+\mathbf{b}_{h}\right) \\
					& \mathbf{h}^{(3)}=g_{1}\left(\mathbf{W}\left(g_{1}\left(\mathbf{W}\left(g_{1}\left(\mathbf{W h}^{(0)}+\mathbf{U} \mathbf{x}^{(1)}+\mathbf{b}_{h}\right)\right)+\mathbf{U} \mathbf{x}^{(2)}+\mathbf{b}_{h}\right)\right)+\mathbf{U} \mathbf{x}^{(3)}+\mathbf{b}_{h}\right) \\
					& \vdots \\
					& \mathbf{h}^{(t)}=g_{1}\left(\mathbf{W h}^{(t-1)}+\mathbf{U} \mathbf{x}^{(t)}+\mathbf{b}_{h}\right) {\color{black} \Leftarrow \text{Salida de la capa recurrente.}}\\
				\end{aligned}
			\end{equation*}}

			{\color{darkgreen}
			\begin{equation*}
				\begin{aligned}
					& \mathbf{y}^{(1)}=g_{2}\left(\mathbf{V h}^{(1)}+\mathbf{b}_{y}\right)=g_{2}\left(\mathbf{V} \cdot g_{1}\left(\mathbf{W h}^{(0)}+\mathbf{U} \mathbf{x}^{(1)}+\mathbf{b}_{h}\right)+\mathbf{b}_{y}\right) \\
					& \mathbf{y}^{(2)}=g_{2}\left(\mathbf{V} \cdot g_{1}\left(\mathbf{W}\left(g_{1}\left(\mathbf{W} \mathbf{h}^{(0)}+\mathbf{U} \mathbf{x}^{(1)}+\mathbf{b}_{h}\right)\right)+\mathbf{U} \mathbf{x}^{(2)}+\mathbf{b}_{h}\right)+\mathbf{b}_{y}\right) \\
					& \mathbf{y}^{(3)}=g_{2}\left(\mathbf{V} \cdot g_{1}\left(\mathbf{W}\left(g_{1}\left(\mathbf{W}\left(g_{1}\left(\mathbf{W h}^{(0)}+\mathbf{U} \mathbf{x}^{(1)}+\mathbf{b}_{h}\right)\right)+\mathbf{U} \mathbf{x}^{(2)}+\mathbf{b}_{h}\right)\right)+\mathbf{U} \mathbf{x}^{(3)}+\mathbf{b}_{h}\right)+\mathbf{b}_{y}\right) \\
					& \vdots \\	
					& \mathbf{y}^{(t)}=g_{2}\left(\mathbf{V h}^{(t)}+\mathbf{b}_{y}\right) {\color{black} \Leftarrow \text{Salida de la red.}}
				\end{aligned}
			\end{equation*}}
		
		\end{column}

		\begin{column}{0.4\textwidth}
			\includegraphics[width=.4\textwidth, right]{imgs/tema4/rnn/MiniModel.pdf}
		\end{column}
	\end{columns}

\end{frame}


\begin{frame}{Despliegue en una RNN}

	Las mismas matrices $\mathbf{U}$, $\mathbf{V}$ y $\mathbf{W}$ se utilizan en distintas etapas de la secuencia de entrada:\\
	\vspace{1em}
	\begin{figure}
		\includegraphics[width=.7\textwidth, center]{imgs/tema4/rnn/DESP1.pdf}
		\label{key}
	\end{figure}

\end{frame}

\begin{frame}{Entrenamiento de las RNN}
	\begin{columns}
		\begin{column}{0.5\textwidth}
			\textbf{Backpropagation through time (BPTT):} 
			\begin{enumerate}
				\item Se presenta una secuencia de etapas con su correspondiente entrada y salida deseada a la red.
				\item Se desarrolla la red y se aplica la propagación hacia delante (etapa a etapa), luego se calculan y acumulan los errores de cada etapa.
				\item Se actualizan los pesos en el sentido contrario al desarrollo de la red (como si se enrollase la red de nuevo).
			\end{enumerate}

		\end{column}

		\begin{column}{0.5\textwidth}
			\includegraphics[width=.9\textwidth]{imgs/tema4/rnn/TRAIN1.pdf}
		\end{column}
	\end{columns}

\end{frame}


\begin{frame}[t]{Arquitecturas de ejemplo: 1 a N}
	\textbf{Generación automática de música:} A partir de una nota de partida se genera una pieza musical, es decir, una secuencia más o menos larga de notas.\\
	\vspace{.3cm}
	\includegraphics[width=.55\textwidth, center]{imgs/tema4/rnn/M1N.pdf}
\end{frame}

\begin{frame}[t]{Arquitecturas de ejemplo: M a 1}
	\textbf{Análisis de sentimientos:} A partir de un texto, predecir si la opinión que refleja es favorable (positiva) o desfavorable (negativa) con respecto a un ítem que se está evaluando.\\
	\vspace{.3cm}
	\includegraphics[width=.55\textwidth, center]{imgs/tema4/rnn/MM1.pdf}
\end{frame}

\begin{frame}[t]{Arquitecturas de ejemplo: M a N}
	\textbf{Traducción automática:} A partir de un texto en un idioma, generar la versión traducida a otro idioma.\\
	\vspace{.3cm}
	\includegraphics[width=.9\textwidth, center]{imgs/tema4/rnn/MMN.pdf}
\end{frame}

\begin{frame}{Exploding/Vanishing Gradients}
	En redes profundas, debido al gran número de productos encadenados, puede producirse los ya descritos:
	\begin{itemize}
		\item \textbf{Exploding gradients:} Crecimiento de gradientes tal que las oscilaciones durante el descenso de gradiente provocan que la red no converja.
		\item \textbf{Vanishing gradients:} Los gradientes se hacen tan pequeños que el descenso apenas es apreciable y la red no aprende.
	\end{itemize}
	\begin{block}{}
		\centering
		¿Puede ser un problema en una red recurrente con una sola capa? $\rightarrow$ {\color{red}\textbf{SI}}
	\end{block}
\end{frame}

\begin{frame}{Exploding/Vanishing Gradients}
	\begin{block}{}
		\centering
		Una secuencia de \textbf{n} elementos hace que la red se comporte como si tuviese \textbf{n} capas
	\end{block}
	\vspace{-1.5em}
	\begin{equation*}
		\scalebox{0.445}{ % Ajusta el factor de escala según tus necesidades
		  $\mathbf{h}^{(10)}=g_1\left(\mathbf{M}\left[g_1\left(\mathbf{M}\left[g_1\left(\mathbf{M}\left[g_1\left(\mathbf{M}\left[g_1\left(\mathbf{M}\left[g_1\left(\mathbf{M}\left[g_1\left(\mathbf{M}\left[g_1\left(\mathbf{M}\left[g_1\left(\mathbf{M}\left[g_1\left(\mathbf{M}\left[h^{(0)}, \mathbf{x}^{(1)}\right]+\mathbf{b}_h\right), \mathbf{x}^{(2)}\right]+\mathbf{b}_h\right), \mathbf{x}^{(3)}\right]+\mathbf{b}_h\right), \mathbf{x}^{(4)}\right]+\mathbf{b}_h\right), \mathbf{x}^{(5)}\right]+\mathbf{b}_h\right), \mathbf{x}^{(6)}\right]+\mathbf{b}_h\right), \mathbf{x}^{(7)}\right]+\mathbf{b}_h\right), \mathbf{x}^{(8)}\right]+\mathbf{b}_h\right), \mathbf{x}^{(9)}\right]+\mathbf{b}_h\right), \mathbf{x}^{(10)}\right]+\mathbf{b}_h\right)$
		}
	\end{equation*}
	\begin{columns}
		\begin{column}{0.45\textwidth}
			\includegraphics[width=.8\textwidth]{imgs/tema4/rnn/DESP1.pdf}
		\end{column}

		\begin{column}{0.55\textwidth}
			Solución:
			\begin{itemize}
				\item \textbf{Exploding:} Estrategias de clipping.
				\item \textbf{Vanishing:} Diseño de nuevos tipos de red:
				\begin{itemize}
					\item Redes GRU.
					\item Redes LSTM.
				\end{itemize}
			\end{itemize}
		\end{column}
	\end{columns}

\end{frame}

\begin{frame}{Vanishing Gradients}
	\begin{itemize}
		\item En redes \textbf{feed forward} este efecto provoca variaciones mínimas en las primeras capas.
		\item En las \textbf{redes recurrentes} implica que la influencia de los elementos de la secuencia de entrada se concentrará en el final de la secuencia.
		\begin{itemize}
			\item \textit{``La {\color{blue}bicicleta}, siendo la alternativa de transporte urbano más ecológica, {\color{blue}es} poco usada''}
			\item \textit{``Las {\color{red}bicicletas}, siendo la alternativa de transporte urbano más ecológica, {\color{red}son} poco usadas''}
		\end{itemize}
	\end{itemize}

	\begin{block}{}
		\centering
		Sería bueno que la red `recordase' que estamos hablando en {\color{blue}singular} o en {\color{red}plural}.
	\end{block}

\end{frame}

\begin{frame}[t]{Referencias}
	\begin{enumerate}
	  \item \textbf{Redes recurrentes (GRU, LSTM y Bidireccionales), Oscar Luaces}
	\end{enumerate}
  \end{frame}

\end{document}
